% ==============================================================================
% DECLARAÇÃO SOBRE O USO DE IA GENERATIVA E TECNOLOGIAS ASSISTIDAS POR IA
% ==============================================================================

\section*{Declaração sobre o uso de IA generativa}
\label{sec:declaracao-uso-ia}

Ferramentas de inteligência artificial generativa foram utilizadas neste
trabalho de quatro formas distintas, descritas também na
\autoref{sec:uso-ia}.

Primeiro, modelos de linguagem foram o \textbf{objeto} do experimento: Claude
Sonnet 5.5 e Claude Haiku 5.5 (Anthropic, via Claude Code 2.1.312), Muse Spark
1.3, DeepSeek v4.1 flash e Beam (via OpenCode v2.2.0) e gpt-oss-20b (OpenAI,
executado localmente com llama.cpp) escreveram o código avaliado.

Segundo, a ferramenta Claude Code, desenvolvida pela Anthropic, com o modelo
Claude Opus 5.5, foi utilizada na construção do orquestrador, dos
\textit{scripts} de análise e dos relatórios intermediários do
\textit{benchmark}, sob direção e revisão do autor. O Muse Spark foi utilizado
como ferramenta na criação de parte do código do repositório-sujeito, anterior
a este trabalho.

Terceiro, os modelos Claude Opus 5.5 (Anthropic), Muse Spark 1.3 e DeepSeek
v4.1 flash atuaram como \textbf{juízes} na terceira camada da análise de
qualidade de código, inclusive sobre o código da própria família, como controle
de autopreferência, sob as salvaguardas descritas na
\autoref{sec:qualidade-metodo}.

Quarto, o texto deste artigo foi redigido com o auxílio do modelo Claude Fable
5.1, desenvolvido pela Anthropic, a partir dos relatórios, dos dados e da
documentação produzidos durante a pesquisa, e foi integralmente revisado pelo
autor. Todo o conteúdo produzido com o auxílio dessas ferramentas foi
analisado, revisado e validado pelo autor, que assume integral responsabilidade
pelo conteúdo desta publicação.
